% Generated draft: requires mathematical and source review.
\documentclass[11pt]{article}
\usepackage[margin=25mm]{geometry}
\usepackage{fontspec}
\setmainfont{Latin Modern Roman}
\usepackage{amsmath,amssymb,mathtools,unicode-math}
\setmathfont{Latin Modern Math}
\usepackage{enumitem,needspace,xurl,hyperref}
\setlength{\parindent}{0pt}
\setlength{\parskip}{5pt}
\setlength{\emergencystretch}{3em}
\newcommand{\N}{\mathbb{N}}
\newcommand{\Z}{\mathbb{Z}}
\newcommand{\Q}{\mathbb{Q}}
\newcommand{\R}{\mathbb{R}}
\newcommand{\C}{\mathbb{C}}
\newcommand{\F}{\mathbb{F}}
\newcommand{\A}{\mathbb{A}}
\newcommand{\PP}{\mathbb{P}}
\newcommand{\E}{\mathbb{E}}
\newcommand{\eps}{\varepsilon}
\newcommand{\dd}{\,\mathrm{d}}
\begin{document}
\section*{UnsolvedMath ID: 30000703}
\textbf{Source problem code:} catalogue entry
\subsection{Definitions and mathematical statement}
Let \(\mathbb D=\{z\in\C:|z|<1\}\), and let \(f:\mathbb D\to\mathbb D\) be holomorphic. Define its normalized hyperbolic derivative magnitude by
\[
 f^{\#}_{\mathbb D}(z)=
 \frac{(1-|z|^2)|f'(z)|}{1-|f(z)|^2}.
\]
The denominator is positive. This is the disc hyperbolic derivative, not the spherical derivative with denominator \(1+|f|^2\). The Schwarz--Pick comparison places it between zero and one.

The source hypothesis is the unrestricted interior limit
\[
 \lim_{\substack{z\to1\\z\in\mathbb D}}f^{\#}_{\mathbb D}(z)=1,
\]
meaning that for every \(\varepsilon>0\) there is \(\delta>0\) such that \(|z-1|<\delta\), \(z\in\mathbb D\), implies \(|f^{\#}_{\mathbb D}(z)-1|<\varepsilon\). All approaches, including tangential ones, are included in this printed limit. It is not replaced by a radial limit, a limsup or convergence along one sequence.

The question asks what boundary behavior of \(f\) at \(1\) is forced by this condition. A precise description might characterize the cluster set
\[
 C(f,1)=\{w\in\overline{\mathbb D}:\exists z_j\in\mathbb D,
 \ z_j\to1,\ f(z_j)\to w\},
\]
or establish specified ordinary/angle limits, angular derivatives or continuation properties. The record does not select a particular conclusion or regularity class, so no finite derivative, continuous extension or automorphism conclusion is inserted.

An angular statement concerns limits within each fixed Stolz region \(|1-z|\leq C(1-|z|)\); it is weaker than an unrestricted approach statement. If an angular value \(\zeta\) exists, its derivative assertion would involve the angular limit of \((f(z)-\zeta)/(z-1)\), not an a priori value of \(f'(1)\). The source does not assume \(f\) is defined at \(1\), that its boundary value is \(1\), or any fixed-point or iteration condition. Asymptotic equality at this boundary point is not the same hypothesis as exact Schwarz--Pick equality at an interior point.

The literature triage checked 2026-08-21 reports boundary Schwarz--Pick and Julia--Carath\'eodory theory giving angular conclusions and extremal rigidity, while saying the standard theorem does not supply a complete unrestricted description. Its partially-solved label and notation-restoration verification remain dated metadata; no particular angular theorem is asserted here without its missing hypotheses.
\subsection{Short English statement}
What value, derivative or other boundary properties follow when the disc hyperbolic derivative magnitude tends to one at \(1\) along every interior approach? The full limit is preserved. It is not replaced by an angular hypothesis or treated as an assumption that \(f(1)=1\) or that \(f\) is an automorphism.
\subsection{Sources}
\begin{itemize}
\item Local UnsolvedMath registry, permanent ID 30000703, source code catalogue entry, ``Boundary Behavior under Asymptotic Schwarz--Pick Equality,'' including the full normalized-derivative limit and open-ended boundary question.
\item \emph{Normal Families and Complex Dynamics (2007)}, Oberwolfach Reports, pp.\ 487--548, DOI 10.4171/owr/2007/09: \url{https://doi.org/10.4171/owr/2007/09}.
\item Supplied reference ``Schwarz lemma at the boundary'': \url{https://arxiv.org/abs/1001.1805}.
\item Supplied reference ``hyperbolic distortion at the boundary'': \url{https://arxiv.org/abs/2410.13965}.
\item Literature triage checked 2026-08-21; local record updated 2026-09-16. No external source has been visited or independently checked here.
\end{itemize}
\end{document}
